\documentclass[aspectratio=169,10pt,english]{beamer}

\input{../../../_preambulo_beamer.tex}
\input{../../../_comandos_beamer.tex}

\selectlanguage{english}

\title{MA1001B - Statistical Modeling for Decision Making}
\subtitle{Lesson 44: One-Way ANOVA: SS, MS, F}
\author{}
\institute{Tecnol\'ogico de Monterrey}
\date{}

\begin{document}

\begin{frame}[plain]
  \vspace{1.2cm}
  {\Large\bfseries MA1001B - Statistical Modeling for Decision Making\par}
  \vspace{0.7cm}
  {\large Lesson 44: One-Way ANOVA: SS, MS, F\par}
  \vspace{0.7cm}
  {\color{TecRojo}\rule{\textwidth}{1.2pt}}
  \vspace{0.8cm}

  MA1001B Faculty\\[0.3cm]
  Term\\[0.3cm]
  Tecnol\'ogico de Monterrey
\end{frame}

\begin{frame}{The Three-Mean Question}
  \begin{alertblock}{Guiding question}
    How can one $F$ ratio test whether three packing-method means differ
    without running an unlabeled collection of $t$ tests?
  \end{alertblock}

  \vspace{0.6em}
  By the end of this lesson, you can:
  \begin{itemize}
    \item read $n$, $\bar y_i$, and $s_i$ for three independent samples;
    \item decompose $SST=SSB+SSE$ and form $F=MSB/MSE$;
    \item compare $F$ with $F_{0.05}(2,33)$ and a $p$-value;
    \item explain why a significant $F$ still does not name a pair.
  \end{itemize}

  \vfill
  \scriptsize All cycle times used today are synthetic. No real company data are used.
\end{frame}

\begin{frame}{Three Independent Samples}
  \small
  \begin{center}
    \begin{tabular}{lrrr}
      \toprule
      \textbf{Method} & $\mathbf{n}$ & \textbf{Mean} & $\mathbf{s}$ \\
      \midrule
      Standard & 12 & 49.547060 & 3.067311 \\
      Guided & 12 & 48.331227 & 2.353358 \\
      Automated & 12 & 43.830670 & 2.660249 \\
      \midrule
      Combined & 36 & 47.236319 &  \\
      \bottomrule
    \end{tabular}
  \end{center}

  \vspace{0.5em}
  \begin{exampleblock}{Hypotheses}
    $H_0$: $\mu_{\text{S}}=\mu_{\text{G}}=\mu_{\text{A}}$.
    $H_1$: at least one mean differs. The factor is packing method; the
    response is cycle time in seconds.
  \end{exampleblock}
\end{frame}

\begin{frame}[fragile]{Step 1: Group Summaries}
  \begin{columns}[T]
    \begin{column}{0.40\textwidth}
      \small
      \begin{lstlisting}
mean_i = y[g == i].mean()
s_i = y[g == i].std(ddof=1)
grand = y.mean()
      \end{lstlisting}

      \begin{block}{Verified output}
        Means $49.547$, $48.331$, $43.831$\\
        $s=3.067$, $2.353$, $2.660$\\
        Grand mean $47.236319$
      \end{block}
    \end{column}
    \begin{column}{0.60\textwidth}
      \centering
      \includegraphics[width=\textwidth]{../figures/one_way_anova_01_group_summaries.png}
      \vspace{0.2em}
      \scriptsize Three-sample boxplots from Step 1
    \end{column}
  \end{columns}
\end{frame}

\begin{frame}{SST, SSB, SSE, and the $F$ Ratio}
  \small
  \[
    SST=SSB+SSE=459.901248,
    \qquad
    F=\frac{MSB}{MSE}=\frac{108.820749}{7.341205}=14.823283.
  \]
  \begin{center}
    \begin{tabular}{lrrrr}
      \toprule
      \textbf{Source} & \textbf{SS} & \textbf{df} & \textbf{MS} & $\mathbf{F}$ \\
      \midrule
      Between & 217.641499 & 2 & 108.820749 & 14.823283 \\
      Error & 242.259749 & 33 & 7.341205 &  \\
      Total & 459.901248 & 35 &  &  \\
      \bottomrule
    \end{tabular}
  \end{center}
\end{frame}

\begin{frame}[fragile]{Step 2: The ANOVA Table}
  \begin{columns}[T]
    \begin{column}{0.40\textwidth}
      \small
      \begin{lstlisting}
MSB = SSB / (k - 1)
MSE = SSE / (N - k)
F = MSB / MSE
      \end{lstlisting}

      \begin{block}{Verified output}
        $SSB=217.641499$\\
        $SSE=242.259749$\\
        Manual $F=14.823283$
      \end{block}
    \end{column}
    \begin{column}{0.60\textwidth}
      \centering
      \includegraphics[width=\textwidth]{../figures/one_way_anova_02_ss_ms_table.png}
      \vspace{0.2em}
      \scriptsize Sum-of-squares decomposition from Step 2
    \end{column}
  \end{columns}
\end{frame}

\begin{frame}{Step 3: A Significant $F$ Does Not Name a Pair}
  \begin{columns}[T]
    \begin{column}{0.42\textwidth}
      \small
      scipy $F=14.823283$, matching the table.

      \vspace{0.4em}
      $p=2.550793\times10^{-5}$,
      $F_{0.05}(2,33)=3.284918$. Reject $H_0$.

      \vspace{0.4em}
      Gaps: $1.215833$, $5.716390$, $4.500557$.

      \vspace{0.5em}
      \textbf{Limit:} a significant $F$ does not say which pair differs.
      Lesson 45 organizes that question as Tukey HSD.
    \end{column}
    \begin{column}{0.58\textwidth}
      \centering
      \includegraphics[width=\textwidth]{../figures/one_way_anova_03_f_test_limit.png}
      \vspace{0.2em}
      \scriptsize Omnibus $F$ and unlabeled gaps from Step 3
    \end{column}
  \end{columns}
\end{frame}

\begin{frame}{Concept Check}
  \begin{enumerate}
    \item Why is $df$ between equal to $2$ in this experiment?
    \item Verify $SSB+SSE=SST$ from the table values.
    \item Why can $F=14.823283$ reject $H_0$ while Standard versus Guided
      remains an open pairwise question?
    \item What does \texttt{scipy.stats.f\_oneway} add after the manual table?
  \end{enumerate}

  \vspace{0.6em}
  \begin{block}{Connection to Lesson 45}
    The session can continue directly into Tukey HSD, familywise error, and
    which packing-method pairs differ.
  \end{block}

  \vfill
  \scriptsize Source: OpenStax, \textit{Introductory Business Statistics 2e},
  Sections 12.2--12.3. CC BY-NC-SA.
\end{frame}

\appendix



\end{document}
